\documentclass[10pt]{article}
\usepackage[utf8]{inputenc}
\usepackage[T1]{fontenc}
\usepackage{amsmath}
\usepackage{amsfonts}
\usepackage{amssymb}
\usepackage[version=4]{mhchem}
\usepackage{stmaryrd}

\title{ECN 711: Government in OG Economy }

\author{Gustavo Ventura}
\date{}


\begin{document}
\maketitle


\section*{Introducing a Government}
\section*{Introducing a Government}
\begin{itemize}
  \item We now a introduce an agent in the economy, a government. A government taxes individuals, consumes goods and can provide for transfers.
  \item We now a introduce an agent in the economy, a government. A government taxes individuals, consumes goods and can provide for transfers.
  \item Let
\end{itemize}

$$
\tau_{t}^{h}=\left[\tau_{t}^{h}(t), \tau_{t}^{h}(t+1)\right]
$$

be the taxes (transfers if negative) that agent $h$ of generation $t$ faces in his/her lifetime.

\begin{itemize}
  \item We now a introduce an agent in the economy, a government. A government taxes individuals, consumes goods and can provide for transfers.
  \item Let
\end{itemize}

$$
\tau_{t}^{h}=\left[\tau_{t}^{h}(t), \tau_{t}^{h}(t+1)\right]
$$

be the taxes (transfers if negative) that agent $h$ of generation $t$ faces in his/her lifetime.

\begin{itemize}
  \item Let $G(t)$ be goods consumed by the government at $t$.
  \item We now a introduce an agent in the economy, a government. A government taxes individuals, consumes goods and can provide for transfers.
  \item Let
\end{itemize}

$$
\tau_{t}^{h}=\left[\tau_{t}^{h}(t), \tau_{t}^{h}(t+1)\right]
$$

be the taxes (transfers if negative) that agent $h$ of generation $t$ faces in his/her lifetime.

\begin{itemize}
  \item Let $G(t)$ be goods consumed by the government at $t$.
  \item Government's budget constraint:
\end{itemize}

$$
G(t)=\sum_{h=1}^{N(t)} \tau_{t}^{h}(t)+\sum_{h=1}^{N(t-1)} \tau_{t-1}^{h}(t)
$$

$40 \times 0 \times 1 \equiv \times, \equiv$ mac

\begin{itemize}
  \item Given taxation, lifetime budget constraint is now:
\end{itemize}

$$
c_{t}^{h}(t)+\frac{c_{t}^{h}(t+1)}{R(t)} \leq w_{t}^{h}(t)-\tau_{t}^{h}(t)+\frac{w_{t}^{h}(t+1)-\tau_{t}^{h}(t+1)}{R(t)}
$$

\begin{itemize}
  \item Given taxation, lifetime budget constraint is now:
\end{itemize}

$$
c_{t}^{h}(t)+\frac{c_{t}^{h}(t+1)}{R(t)} \leq w_{t}^{h}(t)-\tau_{t}^{h}(t)+\frac{w_{t}^{h}(t+1)-\tau_{t}^{h}(t+1)}{R(t)}
$$

\begin{itemize}
  \item Individual decisions are characterized by still characterized by the same first-order conditions. Why?
  \item Given taxation, lifetime budget constraint is now:
\end{itemize}

$$
c_{t}^{h}(t)+\frac{c_{t}^{h}(t+1)}{R(t)} \leq w_{t}^{h}(t)-\tau_{t}^{h}(t)+\frac{w_{t}^{h}(t+1)-\tau_{t}^{h}(t+1)}{R(t)}
$$

\begin{itemize}
  \item Individual decisions are characterized by still characterized by the same first-order conditions. Why?\\
-What is an equilibrium in these circumstances?
\end{itemize}

\section*{Equilibrium Definition}
Let the tax sequence $T$ be

$$
T \equiv\left\{\left(\tau_{t-1}^{h}(t)\right)_{h=1}^{N(t-1)},\left(\tau_{t}^{h}(t)\right)_{h=1}^{N(t)}\right\}_{t=1}^{\infty}
$$

We can now define a competitive equilibrium.

\section*{Equilibrium Definition}
A competitive equilibrium is a tax sequence $\widetilde{T}$, a sequence of government consumption $\{\widetilde{G}(t)\}_{t=1}^{\infty}$, a consumption allocation $\widetilde{C}$, a sequence of lending/borrowing decisions, $\left\{\left(\widetilde{I}^{h}(t)\right)_{h=1}^{N(t)}\right\}_{t=1}^{\infty}$, and interest rates $\{\widetilde{R}(t)\}_{t=1}^{\infty}$ such that\\
(1) Given $\widetilde{T}$ and $\{\widetilde{R}(t)\}_{t=1}^{\infty}$, individuals maximize utility and choose the relevant elements of $\widetilde{C}$.\\
(2) Market clearing, i.e. for all $t \geq 1$

$$
\sum_{h=1}^{N(t)} \pi^{h}(t)=0
$$

(3) The government budget balances for all $t \geq 1$

$$
\sum_{h=1}^{N(t)} \widetilde{\tau}_{t}^{h}(t)+\sum_{h=1}^{N(t-1)} \widetilde{\tau}_{t-1}^{h}(t)=\widetilde{G}(t)
$$

\section*{Remarks}
\section*{Remarks}
\begin{itemize}
  \item New condition: budget balance for the government!
\end{itemize}

\section*{Remarks}
\begin{itemize}
  \item New condition: budget balance for the government!\\
-What is the feasibility condition now?
\end{itemize}

\section*{Remarks}
\begin{itemize}
  \item New condition: budget balance for the government!
  \item What is the feasibility condition now?
  \item Answer:
\end{itemize}

$$
C(t)+G(t)=Y(t)
$$

\begin{itemize}
  \item New condition: budget balance for the government!
  \item What is the feasibility condition now?
  \item Answer:
\end{itemize}

$$
C(t)+G(t)=Y(t)
$$

\begin{itemize}
  \item Exercise: show that an allocation $\widetilde{C}$ and a path of government consumption are feasible if and only if the market for consumption loans clears in all dates.
\end{itemize}

Exercise: intergenerational transfers

\section*{Exercise: intergenerational transfers}
\begin{itemize}
  \item Identical individuals within a cohort of size $N$. Preferences:
\end{itemize}

$$
u\left(c_{t}(t), c_{t}(t+1)=\log \left(c_{t}(t)\right)+\log \left(c_{t}(t+1)\right)\right.
$$

\section*{Exercise: intergenerational transfers}
\begin{itemize}
  \item Identical individuals within a cohort of size $N$. Preferences:
\end{itemize}

$$
u\left(c_{t}(t), c_{t}(t+1)=\log \left(c_{t}(t)\right)+\log \left(c_{t}(t+1)\right)\right.
$$

\begin{itemize}
  \item Endowments: $\omega_{y}$ and $\omega_{o}$.
\end{itemize}

\section*{Exercise: intergenerational transfers}
\begin{itemize}
  \item Identical individuals within a cohort of size $N$. Preferences:
\end{itemize}

$$
u\left(c_{t}(t), c_{t}(t+1)=\log \left(c_{t}(t)\right)+\log \left(c_{t}(t+1)\right)\right.
$$

\begin{itemize}
  \item Endowments: $\omega_{y}$ and $\omega_{o}$.
  \item Government taxes each young with tax $\tau$, to pay for transfer $b$ to each old. No government consumption.
\end{itemize}

\section*{Exercise: intergenerational transfers}
\begin{itemize}
  \item Identical individuals within a cohort of size $N$. Preferences:
\end{itemize}

$$
u\left(c_{t}(t), c_{t}(t+1)=\log \left(c_{t}(t)\right)+\log \left(c_{t}(t+1)\right)\right.
$$

\begin{itemize}
  \item Endowments: $\omega_{y}$ and $\omega_{o}$.
  \item Government taxes each young with tax $\tau$, to pay for transfer $b$ to each old. No government consumption.
  \item How are equilibrium interest rates and allocations affected?
\end{itemize}

Exercise: intergenerational transfers

\section*{Exercise: intergenerational transfers}
\begin{itemize}
  \item Optimal borrowing/lending decision:
\end{itemize}

$$
I(t)=\frac{\omega_{y}-\tau}{2}-\frac{\omega_{o}+b}{2 R(t)} .
$$

\section*{Exercise: intergenerational transfers}
\begin{itemize}
  \item Optimal borrowing/lending decision:
\end{itemize}

$$
I(t)=\frac{\omega_{y}-\tau}{2}-\frac{\omega_{o}+b}{2 R(t)}
$$

\begin{itemize}
  \item In equilibrium, $\widetilde{\tau}=\widetilde{b}$. Thus, equilibrium interest rate is
\end{itemize}

$$
\widetilde{R}(t)=\widetilde{R}=\frac{\omega_{o}+\tau}{\omega_{y}-\tau}
$$

\section*{Exercise: intergenerational transfers}
\begin{itemize}
  \item Optimal borrowing/lending decision:
\end{itemize}

$$
I(t)=\frac{\omega_{y}-\tau}{2}-\frac{\omega_{o}+b}{2 R(t)}
$$

\begin{itemize}
  \item In equilibrium, $\widetilde{\tau}=\widetilde{b}$. Thus, equilibrium interest rate is
\end{itemize}

$$
\widetilde{R}(t)=\widetilde{R}=\frac{\omega_{o}+\tau}{\omega_{y}-\tau}
$$

\begin{itemize}
  \item Equilibrium interest rate is higher than in the absence of taxes/transfers. Why? What are the implications of population growth?
\end{itemize}

Fluctuating government consumption

\section*{Fluctuating government consumption}
\begin{itemize}
  \item Identical individuals within a cohort of size $N$. Preferences:
\end{itemize}

$$
u\left(c_{t}(t), c_{t}(t+1)=\log \left(c_{t}(t)\right)+\log \left(c_{t}(t+1)\right) .\right.
$$

\section*{Fluctuating government consumption}
\begin{itemize}
  \item Identical individuals within a cohort of size $N$. Preferences:
\end{itemize}

$$
u\left(c_{t}(t), c_{t}(t+1)=\log \left(c_{t}(t)\right)+\log \left(c_{t}(t+1)\right)\right.
$$

\begin{itemize}
  \item Endowments: $\omega_{y}$ and $\omega_{o}$.
\end{itemize}

\section*{Fluctuating government consumption}
\begin{itemize}
  \item Identical individuals within a cohort of size $N$. Preferences:
\end{itemize}

$$
u\left(c_{t}(t), c_{t}(t+1)=\log \left(c_{t}(t)\right)+\log \left(c_{t}(t+1)\right)\right.
$$

\begin{itemize}
  \item Endowments: $\omega_{y}$ and $\omega_{o}$.
  \item Only young pay taxes, $\tau_{t}(t)$. Government consumption fluctuates: $G(t)=G_{H}$ if $t$ is even. $G(t)=G_{L}$ if $t$ is odd. $G_{H}>G_{L}$.
\end{itemize}

Fluctuating government consumption

\section*{Fluctuating government consumption}
\begin{itemize}
  \item Let $g(t)=G(t) / N$, or government consumption per young person. In equilibrium, $\widetilde{g(t)}=\widetilde{\tau_{t}(t)}$
\end{itemize}

\section*{Fluctuating government consumption}
\begin{itemize}
  \item Let $g(t)=G(t) / N$, or government consumption per young person. In equilibrium, $\widetilde{g(t)}=\widetilde{\tau_{t}(t)}$
  \item It follows that the equilibrium path of interest rates obeys\\
$\widetilde{R}(t)=\frac{w_{o}}{w_{y}-g_{H}}$ if $t$ is even.\\
$\widetilde{R}(t)=\frac{w_{o}}{w_{y}-g_{L}}$ if $t$ is odd.
\end{itemize}

\section*{Fluctuating government consumption}
\begin{itemize}
  \item Let $g(t)=G(t) / N$, or government consumption per young person. In equilibrium, $\widetilde{g(t)}=\widetilde{\tau_{t}(t)}$
  \item It follows that the equilibrium path of interest rates obeys\\
$\widetilde{R}(t)=\frac{w_{o}}{w_{y}-g_{H}}$ if $t$ is even.\\
$\widetilde{R}(t)=\frac{w_{o}}{w_{y}-g_{L}}$ if $t$ is odd.
  \item Why are equilibrium interest rates higher when government consumption is high?
\end{itemize}

\section*{Public Debt}
\begin{itemize}
  \item We now assume that the government can issue debt in each date, $B(t)$, at the price $p(t)$. The budget constraint of the government is now:
\end{itemize}

$$
\sum_{h=1}^{N(t)} \tau_{t}^{h}(t)+\sum_{h=1}^{N(t-1)} \tau_{t-1}^{h}(t)+p(t) B(t)=G(t)+B(t-1)
$$

\begin{itemize}
  \item We now assume that the government can issue debt in each date, $B(t)$, at the price $p(t)$. The budget constraint of the government is now:
\end{itemize}

$$
\sum_{h=1}^{N(t)} \tau_{t}^{h}(t)+\sum_{h=1}^{N(t-1)} \tau_{t-1}^{h}(t)+p(t) B(t)=G(t)+B(t-1)
$$

\begin{itemize}
  \item Individual budget constraints are now:
\end{itemize}


\begin{equation*}
c_{t}^{h}(t)=w_{t}^{h}(t)-\tau_{t}^{h}(t)-l^{h}(t)-p(t) b^{h}(t), \tag{1}
\end{equation*}



\begin{equation*}
c_{t}^{h}(t+1)=w_{t}^{h}(t+1)-\tau_{t}^{h}(t+1)+R(t) l^{h}(t)+b^{h}(t), \tag{2}
\end{equation*}


\section*{Public Debt}
\begin{itemize}
  \item (1) and (2) imply lifetime constraint
\end{itemize}

$$
\begin{aligned}
c_{t}^{h}(t)+\frac{c_{t}^{h}(t+1)}{R(t)} & =w_{t}^{h}(t)-\tau_{t}^{h}(t)+\frac{w_{t}^{h}(t+1)-\tau_{t}^{h}(t+1)}{R(t)} \\
& -b^{h}(t)\left[p(t)-\frac{1}{R(t)}\right]
\end{aligned}
$$

\begin{itemize}
  \item (1) and (2) imply lifetime constraint
\end{itemize}

$$
\begin{aligned}
c_{t}^{h}(t)+\frac{c_{t}^{h}(t+1)}{R(t)} & =w_{t}^{h}(t)-\tau_{t}^{h}(t)+\frac{w_{t}^{h}(t+1)-\tau_{t}^{h}(t+1)}{R(t)} \\
& -b^{h}(t)\left[p(t)-\frac{1}{R(t)}\right] .
\end{aligned}
$$

\begin{itemize}
  \item Note that in equilibrium, it has to be that $p(t)=1 / R(t)$. Why? Lifetime budget constraint is the same as before (last term disappears!).
\end{itemize}

\section*{Equilibrium with Public Debt}
\begin{itemize}
  \item Hence, individuals are indifferent between holding debt or consumption loans. Let
\end{itemize}

$$
s^{h}(t) \equiv l^{h}(t)+p(t) b^{h}(t)
$$

denote what they are willing to lend (or borrow) to the government or others.

\section*{Equilibrium with Public Debt}
\begin{itemize}
  \item Hence, individuals are indifferent between holding debt or consumption loans. Let
\end{itemize}

$$
s^{h}(t) \equiv l^{h}(t)+p(t) b^{h}(t)
$$

denote what they are willing to lend (or borrow) to the government or others.

\begin{itemize}
  \item A definition of equilibrium includes now the path of government debt, $\{\widetilde{B}(t)\}_{t=1}^{\infty}$.
\end{itemize}

\section*{Equilibrium with Public Debt}
\begin{itemize}
  \item Hence, individuals are indifferent between holding debt or consumption loans. Let
\end{itemize}

$$
s^{h}(t) \equiv l^{h}(t)+p(t) b^{h}(t)
$$

denote what they are willing to lend (or borrow) to the government or others.

\begin{itemize}
  \item A definition of equilibrium includes now the path of government debt, $\{\widetilde{B}(t)\}_{t=1}^{\infty}$.
  \item Define an equilibrium with public debt. What is the market clearing condition now?
\end{itemize}

\section*{Equilibrium with Public Debt}
\begin{itemize}
  \item Hence, individuals are indifferent between holding debt or consumption loans. Let
\end{itemize}

$$
s^{h}(t) \equiv l^{h}(t)+p(t) b^{h}(t)
$$

denote what they are willing to lend (or borrow) to the government or others.

\begin{itemize}
  \item A definition of equilibrium includes now the path of government debt, $\{\widetilde{B}(t)\}_{t=1}^{\infty}$.
  \item Define an equilibrium with public debt. What is the market clearing condition now?
\end{itemize}

\section*{Equilibrium with Public Debt}
\begin{itemize}
  \item Market Clearing condition now:
\end{itemize}

$$
\sum_{h=1}^{N(t)} \widetilde{s}^{h}(t)=\underbrace{\sum_{h=1}^{N(t)} \widetilde{l}^{h}(t)+\sum_{h=1}^{N(t)} \frac{\widetilde{b}^{h}(t)}{\widetilde{R}(t)}}_{\text {borrowing/lending plus debt purchases }}=\frac{\widetilde{B}(t)}{\widetilde{R}(t)} .
$$

\section*{Equilibrium with Public Debt}
\begin{itemize}
  \item Market Clearing condition now:
\end{itemize}

$$
\sum_{h=1}^{N(t)} \widetilde{s}^{h}(t)=\underbrace{\sum_{h=1}^{N(t)} \widetilde{I}^{h}(t)+\sum_{h=1}^{N(t)} \frac{\widetilde{b}^{h}(t)}{\widetilde{R}(t)}}_{\text {borrowing/lending plus debt purchases }}=\frac{\widetilde{B}(t)}{\widetilde{R}(t)} .
$$

\begin{itemize}
  \item Remark: Note that there is no net borrowing and lending in any equilibrium. Why?
\end{itemize}

Exercise: transfers with debt

\section*{Exercise: transfers with debt}
\begin{itemize}
  \item Identical individuals within a cohort of size $N$. Preferences:
\end{itemize}

$$
u\left(c_{t}(t), c_{t}(t+1)=\log \left(c_{t}(t)\right)+\log \left(c_{t}(t+1)\right) .\right.
$$

\section*{Exercise: transfers with debt}
\begin{itemize}
  \item Identical individuals within a cohort of size $N$. Preferences:
\end{itemize}

$$
u\left(c_{t}(t), c_{t}(t+1)=\log \left(c_{t}(t)\right)+\log \left(c_{t}(t+1)\right) .\right.
$$

\begin{itemize}
  \item Endowments: $\omega_{y}=2$ and $\omega_{o}=1$.
\end{itemize}

\section*{Exercise: transfers with debt}
\begin{itemize}
  \item Identical individuals within a cohort of size $N$. Preferences:
\end{itemize}

$$
u\left(c_{t}(t), c_{t}(t+1)=\log \left(c_{t}(t)\right)+\log \left(c_{t}(t+1)\right) .\right.
$$

\begin{itemize}
  \item Endowments: $\omega_{y}=2$ and $\omega_{o}=1$.
  \item $G(t)=0$ all $t$. Unexpectedly, the government transfers $\triangle=1 / 2$ to current old at $t_{0}$, financed with debt, to be fully repaid by the young at $t_{0}+1$.
\end{itemize}

\section*{Exercise: transfers with debt}
\begin{itemize}
  \item Identical individuals within a cohort of size $N$. Preferences:
\end{itemize}

$$
u\left(c_{t}(t), c_{t}(t+1)=\log \left(c_{t}(t)\right)+\log \left(c_{t}(t+1)\right) .\right.
$$

\begin{itemize}
  \item Endowments: $\omega_{y}=2$ and $\omega_{o}=1$.
  \item $G(t)=0$ all $t$. Unexpectedly, the government transfers $\triangle=1 / 2$ to current old at $t_{0}$, financed with debt, to be fully repaid by the young at $t_{0}+1$.
  \item Obtain the path of interest rates in this economy. Show that $R(t)=1 / 2$ if $t<t_{0}$ or $t>t_{0}+1$. Further, $R\left(t_{0}\right)=1$ and $R\left(t_{0}+1\right)=2 / 3$.
\end{itemize}

\end{document}